\frfilename{mt113.tex}
\versiondate{6.4.05}
\copyrightdate{1999}

\def\chaptername{Ruang ukur}
\def\sectionname{Ukuran luar dan konstruksi \Caratheodory}

\newsection{113}

Saya memperkenalkan metode terpenting untuk mengonstruksi ukuran.

\leader{113A}{Ukuran \dvrocolon{luar}}\cmmnt{ Sekarang saya sampai pada
definisi dasar ketiga dalam bab ini.

\medskip

\noindent}{\bf Definisi} Misalkan $X$ suatu himpunan. Sebuah {\bf ukuran luar} pada
$X$ adalah fungsi $\theta:\Cal PX\to[0,\infty]$ sedemikian sehingga

\qquad(i) $\theta \emptyset=0$,

\qquad(ii) jika $A\subseteq B\subseteq X$ maka $\theta A\le\theta B$,

\qquad(iii) untuk setiap barisan $\langle A_n\rangle_{n\in\Bbb N}$ dari
subhimpunan-subhimpunan $X$,
$\theta(\bigcup_{n\in\Bbb N}A_n)\le\sum_{n=0}^{\infty}\theta A_n$.

\cmmnt{
\leader{113B}{Catatan (a)} Untuk komentar tentang penggunaan
`$\infty$', lihat 112B.

\header{113Bb}{\bf (b)} Sekali lagi, ukuran-ukuran luar terpenting
harus menunggu hingga
\S\S114-115. Gagasan ukuran `luar' dari suatu himpunan $A$ ialah
bahwa ukuran itu seharusnya menjadi semacam batas atas bagi kemungkinan
nilai ukuran $A$. Jika kita beruntung, nilai ini mungkin benar-benar
merupakan ukuran $A$; tetapi hal itu agaknya hanya berlaku bagi himpunan
dengan batas yang cukup mulus.

\header{113Bc}{\bf (c)} Dengan menggabungkan (i) dan (iii) dari definisi
tersebut, kita melihat bahwa
jika $\theta$ adalah ukuran luar pada $X$, dan $A$, $B$ adalah dua subhimpunan
$X$, maka $\theta(A\cup B)\le\theta A+\theta B$; bandingkan 112Ca dan 112Cc.
}%end of comment

\leader{113C}{Metode \Caratheodory: Teorema} Misalkan $X$ suatu himpunan dan
$\theta$ suatu ukuran luar pada $X$. Tetapkan

\Centerline{$\Sigma=\{E:E\subseteq X,\,
\theta A=\theta(A\cap E)+\theta(A\setminus E)$ untuk setiap $A\subseteq X\}$.}

\noindent Maka $\Sigma$ adalah aljabar-$\sigma$ atas subhimpunan-subhimpunan $X$.
Definisikan $\mu:\Sigma\to[0,\infty]$ dengan menetapkan $\mu E=\theta E$ untuk
$E\in\Sigma$; maka $(X,\Sigma,\mu)$ adalah ruang ukur.

\proof{{\bf (a)} Langkah pertama ialah mengamati bahwa untuk setiap $E$,
$A\subseteq X$ berlaku
$\theta(A\cap E)+\theta(A\setminus E)\ge\theta A$, berdasarkan 113Bc; sehingga

\Centerline{$\Sigma=\{E:E\subseteq X,\,
\theta A\ge\theta(A\cap E)+\theta(A\setminus E)$
untuk setiap $A\subseteq X\}$.}


\medskip

{\bf (b)} Jelas bahwa $\emptyset\in\Sigma$, karena

\Centerline{$\theta(A\cap\emptyset)+\theta(A\setminus\emptyset)
=\theta\emptyset+\theta A=\theta A$}

\noindent untuk setiap $A\subseteq X$. Jika $E\in\Sigma$, maka
$X\setminus E\in\Sigma$, karena

\Centerline{$\theta(A\cap(X\setminus E))+\theta(A\setminus(X\setminus E))
=\theta(A\setminus E)+\theta(A\cap E)=\theta A$}

\noindent untuk setiap $A\subseteq X$.

\medskip

{\bf (c)} Sekarang misalkan $E$, $F\in\Sigma$ dan $A\subseteq X$. Maka

\ifdim\pagewidth<391pt
\startsideshiftedpicture
\hskip-20pt
\sideshiftedpicture{mt113c1}{-15pt}{90.05pt}{100.68pt}
\sideshiftedpicture{mt113c2}{-15pt}{90.05pt}{100.68pt}
\sideshiftedpicture{mt113c3}{-15pt}{90.05pt}{100.68pt}
\sideshiftedpicture{mt113c4}{-15pt}{90.05pt}{100.68pt}
\else
\startsideshiftedpicture
\hskip10pt
\sideshiftedpicture{mt113c1}{0pt}{97.05pt}{100.68pt}
\sideshiftedpicture{mt113c2}{0pt}{97.05pt}{100.68pt}
\sideshiftedpicture{mt113c3}{0pt}{97.05pt}{100.68pt}
\sideshiftedpicture{mt113c4}{0pt}{97.05pt}{100.68pt}
\fi

\ifdim\pagewidth=390pt\vskip-15pt\fi
\ifdim\pagewidth>467pt\vskip-10pt\fi

$$\eqalignno{\theta(A\cap(&E\cup F))+\theta(A\setminus(E\cup F))
    &\text{diagram (i)}\cr
&=\theta(A\cap(E\cup F)\cap E)+\theta(A\cap(E\cup F)\setminus E)
  +\theta(A\setminus(E\cup F))
    &\text{diag.\ (ii)}\cr
\noalign{\noindent (karena $E\in\Sigma$ dan
$A\cap(E\cup F)\subseteq X$)}
&=\theta(A\cap E)+\theta((A\setminus E)\cap F)
  +\theta((A\setminus E)\setminus F)\cr
&=\theta(A\cap E)+\theta(A\setminus E)
    &\text{diag.\ (iii)}\cr
\noalign{\noindent (karena $F\in\Sigma$)}
&=\theta A
    &\text{diag.\ (iv)}\cr}$$

\noindent (sekali lagi karena $E\in\Sigma$). Karena $A$ sebarang,
$E\cup F\in\Sigma$.

\medskip

{\bf (d)} Dengan demikian $\Sigma$ tertutup terhadap gabungan dua himpunan dan komplemen,
serta memuat $\emptyset$.
%Observe that it follows that
%$E\setminus F=X\setminus(F\cup(X\setminus E))$ belongs to $\Sigma$ whenever $E$,
%$F\in\Sigma$.
Sekarang misalkan $\langle E_n\rangle_{n\in\Bbb N}$ adalah
barisan di dalam $\Sigma$, dengan $E=\bigcup_{n\in\Bbb N}E_n$. Tetapkan

\Centerline{$G_n=\bigcup_{m\le n}E_m$;}

\noindent maka $G_n\in\Sigma$ untuk setiap $n$, melalui induksi terhadap $n$. Tetapkan

\Centerline{$F_0=G_0=E_0$,\quad
$F_n=G_n\setminus G_{n-1}=E_n\setminus G_{n-1}$ untuk $n\ge 1$;}

\noindent maka
%every $F_n$ belongs to $\Sigma$, and
$E=\bigcup_{n\in\Bbb N}F_n=\bigcup_{n\in\Bbb N}G_n$.

Ambil sembarang $n\ge 1$ dan sembarang $A\subseteq X$. Maka

$$\eqalign{\theta(A\cap G_n)
&=\theta(A\cap G_n\cap G_{n-1})+\theta(A\cap G_n\setminus G_{n-1})\cr
&=\theta(A\cap G_{n-1})+\theta(A\cap F_n).\cr}$$

\noindent Induksi terhadap $n$ menunjukkan bahwa $\theta(A\cap
G_n)=\sum_{m=0}^n\theta(A\cap F_m)$ untuk setiap $n\ge 0$.

Misalkan $A\subseteq X$. Maka $A\cap
E=\bigcup_{n\in\Bbb N}A\cap F_n$, sehingga

$$\eqalign{\theta(A\cap E)
&\le\sum_{n=0}^{\infty}\theta(A\cap F_n)\cr
&=\lim_{n\to\infty}\sum_{m=0}^n\theta(A\cap F_m)
=\lim_{n\to\infty}\theta(A\cap G_n).}$$

\noindent Di pihak lain,

$$\eqalign{\theta(A\setminus E)
&=\theta(A\setminus\bigcup_{n\in\Bbb N}G_n)\cr
&\le\inf_{n\in\Bbb N}\theta(A\setminus G_n)
=\lim_{n\to\infty}\theta(A\setminus G_n),}$$

\noindent dengan menggunakan 113A(ii) untuk melihat bahwa
$\langle\theta(A\setminus G_n)\rangle_{n\in\Bbb N}$ bersifat tak-menaik
dan bahwa $\theta(A\setminus E)\le \theta(A\setminus G_n)$ untuk setiap $n$.
Oleh karena itu

$$\eqalign{\theta(A\cap E)+\theta(A\setminus E)
&\le\lim_{n\to\infty}\theta(A\cap G_n)
+\lim_{n\to\infty}\theta(A\setminus G_n)\cr
&=\lim_{n\to\infty}(\theta(A\cap G_n)+\theta(A\setminus G_n))
=\theta A\cr}$$

\noindent karena setiap $G_n$ termasuk dalam $\Sigma$, sehingga $\theta(A\cap
G_n)+\theta(A\setminus G_n)=\theta A$ untuk setiap $n$. Namun, karena $A$
sebarang, $E\in\Sigma$, menurut catatan dalam (a) di atas.

Karena $\langle E_n\rangle_{n\in\Bbb N}$ sebarang,
syarat (iii) dari 111A terpenuhi, dan $\Sigma$ adalah aljabar-$\sigma$
atas subhimpunan-subhimpunan $X$.

\medskip

{\bf (e)} Sekarang mari kita beralih ke $\mu$, yakni pembatasan $\theta$ pada
$\Sigma$, dan
Definisi 112A. Tentu saja $\mu\emptyset=\theta\emptyset=0$. Jadi, misalkan
$\langle E_n\rangle_{n\in\Bbb N}$ sebarang barisan saling lepas di dalam $\Sigma$.
Tetapkan $G_n=\bigcup_{m\le n}E_m$ untuk setiap $n$, seperti dalam (d), dan

\Centerline{$E=\bigcup_{n\in\Bbb N}E_n=\bigcup_{n\in\Bbb N}G_n$.}

\noindent Seperti dalam (d),

$$\eqalign{\mu G_{n+1}
=\theta G_{n+1}
&=\theta(G_{n+1}\cap E_{n+1})+\theta(G_{n+1}\setminus E_{n+1})\cr
&=\theta E_{n+1}+\theta G_n
=\mu E_{n+1}+\mu G_n\cr}$$

\noindent untuk setiap $n$, sehingga $\mu G_n=\sum_{m=0}^n\mu E_m$ untuk setiap $n$.

Sekarang

\Centerline{$\mu E
=\theta E
\le\sum_{n=0}^{\infty}\theta E_n
=\sum_{n=0}^{\infty}\mu E_n$.}

\noindent Tetapi juga

\Centerline{$\mu E=\theta E\ge\theta G_n=\mu G_n=\sum_{m=0}^n\mu E_m$}

\noindent untuk setiap $n$, sehingga $\mu E\ge\sum_{n=0}^{\infty}\mu E_n$.

Oleh karena itu $\mu E=\sum_{n=0}^{\infty}\mu E_n$. Karena
$\langle E_n\rangle_{n\in\Bbb N}$ sebarang, 112A(iii-$\beta$)
terpenuhi dan $(X,\Sigma,\mu)$ adalah ruang ukur.
}%end of proof of 113C

\leader{113D}{Catatan} Perhatikan\cmmnt{ dari (a) dalam bukti di atas} bahwa dalam
konstruksi ini

\cmmnt{\Centerline{$\Sigma=\{E:E\subseteq X,\,\theta(A\cap
E)+\theta(A\setminus
E)\le\theta A$ untuk setiap $A\subseteq X\}$\dvro{,}{.}}

\noindent Karena $\theta(A\cap E)+\theta(A\setminus E)$
pasti kurang dari atau sama dengan $\theta A$ ketika $\theta A=\infty$,}

\Centerline{$\Sigma=\{E:E\subseteq X,\,\theta(A\cap E)+\theta(A\setminus
E)\le\theta A$ apabila $A\subseteq X$ dan $\theta A<\infty\}$.}

\exercises{
\leader{113X}{Latihan dasar $\pmb{>}$(a)}
%\spheader 113Xa
Misalkan $X$ suatu himpunan dan $\theta$ suatu ukuran luar pada $X$, dan
misalkan $\mu$ ukuran pada $X$ yang didefinisikan dari $\theta$ dengan
metode \Caratheodory. Tunjukkan bahwa jika $\theta A=0$, maka $\mu$
mengukur $A$, sehingga suatu himpunan $A\subseteq X$ terabaikan-$\mu$
jika dan hanya jika $\theta A=0$, dan $\mu$ bersifat `lengkap' dalam
arti 112Df.

\spheader 113Xb Misalkan $X$ suatu himpunan.
{(i)} Tunjukkan bahwa jika $\theta_1$, $\theta_2$ adalah ukuran luar pada
$X$, maka $\theta_1+\theta_2$ juga merupakan ukuran luar, dengan menetapkan
$(\theta_1+\theta_2)(A)=\theta_1A+\theta_2A$ untuk setiap $A\subseteq X$.
{(ii)} Tunjukkan bahwa jika
$\langle\theta_i\rangle_{i\in I}$ adalah sembarang keluarga tak kosong yang
terdiri atas ukuran-ukuran luar pada $X$, maka $\theta=\sup_{i\in I}\theta_i$ juga merupakan
ukuran luar, dengan menetapkan $\theta A=\sup_{i\in I}\theta_iA$ untuk
setiap $A\subseteq X$. {(iii)}
Tunjukkan bahwa jika $\theta_1$, $\theta_2$ adalah ukuran luar pada $X$,
maka $\theta_1\wedge\theta_2$ juga merupakan ukuran luar, dengan menetapkan

\Centerline{$(\theta_1\wedge\theta_2)(A)
=\inf\{\theta_1B+\theta_2(A\setminus B):B\subseteq A\}$}

\noindent untuk setiap $A\subseteq X$.

\sqheader 113Xc Misalkan $X$ dan $Y$ himpunan, $\theta$ suatu
ukuran luar pada $X$, dan $f:X\to Y$ suatu fungsi. Tunjukkan bahwa fungsional
$B\mapsto\theta(f^{-1}[B]):\Cal PY\to[0,\infty]$ adalah ukuran luar pada
$Y$.

\sqheader 113Xd Misalkan $X$ suatu himpunan dan $\theta$ suatu ukuran luar
pada $X$; misalkan $Y$ sembarang subhimpunan dari $X$. {(i)} Tunjukkan bahwa
$\theta\restrp\Cal PY$, pembatasan $\theta$ pada subhimpunan-subhimpunan
$Y$, adalah ukuran luar pada $Y$. {(ii)} Tunjukkan bahwa jika $E\subseteq X$
diukur oleh ukuran pada $X$ yang didefinisikan dari $\theta$ dengan
metode \Caratheodory, maka $E\cap Y$ diukur oleh ukuran pada
$Y$ yang didefinisikan dari $\theta\restrp\Cal PY$.

\sqheader 113Xe Misalkan $X$ dan $Y$ himpunan, $\theta$ suatu
ukuran luar pada $Y$, dan $f:X\to Y$ suatu fungsi. Tunjukkan bahwa fungsional
$A\mapsto\theta(f[A]):\Cal PX\to[0,\infty]$ adalah ukuran luar.

\spheader 113Xf Misalkan $X$ dan $Y$ himpunan, $\theta$ suatu
ukuran luar pada $X$, dan $R\subseteq X\times Y$ suatu relasi. Tunjukkan bahwa
pemetaan $B\mapsto\theta(R^{-1}[B]):\Cal PY\to[0,\infty]$ adalah ukuran luar
pada $Y$, dengan $R^{-1}[B]=\{x:\exists\,y\in B,\,(x,y)\in R\}$
(1A1Bc). Jelaskan bagaimana hasil ini merupakan perumuman bersama dari (c),
(d-i), dan (e) di atas, serta bagaimana hasil itu dapat dibuktikan dengan menggabungkannya.

\spheader 113Xg Misalkan $X$ suatu himpunan dan $\theta$ suatu ukuran luar
pada $X$. Andaikan $E\subseteq X$ diukur oleh ukuran pada
$X$ yang didefinisikan dari $\theta$ dengan metode \Caratheodory. Tunjukkan bahwa
$\theta(E\cap A)+\theta(E\cup A)=\theta E+\theta A$ untuk setiap
$A\subseteq X$.

\spheader 113Xh Misalkan $X$ suatu himpunan dan $\theta:\Cal PX\to[0,\infty]$ suatu fungsional sedemikian sehingga $\theta\emptyset=0$, $\theta A\le\theta B$ setiap kali $A\subseteq B\subseteq X$, dan $\theta(A\cup B)\le\theta A+\theta B$ setiap kali $A$, $B\subseteq X$. Tetapkan

\Centerline{$\Sigma=\{E:E\subseteq X,\,
\theta A=\theta(A\cap E)+\theta(A\setminus E)$ untuk setiap
$A\subseteq X\}$.}

\noindent Tunjukkan bahwa $\emptyset$, $X\setminus E$, dan $E\cup F$ termasuk dalam
$\Sigma$ untuk semua $E$,
$F\in\Sigma$, sehingga $E\setminus F$, $E\cap F\in\Sigma$ untuk semua $E$,
$F\in\Sigma$. Tunjukkan bahwa $\theta(E\cup F)=\theta E+\theta F$ setiap kali
$E$, $F\in\Sigma$ dan $E\cap F=\emptyset$.

\leader{113Y}{Latihan lanjutan (a)}
%\spheader 113Ya
Misalkan $(X,\Sigma,\mu)$ suatu ruang ukur. Untuk $A\subseteq X$
tetapkan $\mu^*A=\inf\{\mu E:E\in\Sigma,\,A\subseteq E\}$. Tunjukkan bahwa
untuk setiap $A\subseteq X$ infimum itu dicapai, yakni, terdapat
$E\in\Sigma$ sedemikian sehingga $A\subseteq E$ dan $\mu E=\mu^*A$. Tunjukkan bahwa
$\mu^*$ adalah ukuran luar pada $X$.

\spheader 113Yb Misalkan $(X,\Sigma,\mu)$ suatu ruang ukur dan $D$
sembarang subhimpunan dari
$X$. Tunjukkan bahwa $\Sigma_D=\{E\cap D:E\in\Sigma\}$ adalah aljabar-$\sigma$
atas subhimpunan-subhimpunan $D$. Tetapkan $\mu_D=\mu^*\restr \Sigma_D$, yakni fungsi dengan
domain $\Sigma_D$ sedemikian sehingga $\mu_DB=\mu^*B$ untuk setiap $B\in\Sigma_D$,
dengan $\mu^*$ didefinisikan seperti dalam (a) di atas; tunjukkan bahwa
$(D,\Sigma_D,\mu_D)$ adalah suatu ruang ukur. ($\mu_D$ adalah {\bf ukuran
subruang} pada $D$.)

\spheader 113Yc Misalkan $(X,\Sigma,\mu)$ suatu ruang ukur dan misalkan
$\mu^*$ adalah
ukuran luar terkait pada $X$, seperti dalam 113Ya. Misalkan $\check\mu$ adalah
ukuran pada $X$ yang dikonstruksi dengan metode \Caratheodory\ dari $\mu^*$,
dengan $\check\Sigma$ sebagai domainnya. Tunjukkan bahwa
$\Sigma\subseteq\check\Sigma$ dan
bahwa $\check\mu$ memperluas $\mu$.

\spheader 113Yd Misalkan $X$ suatu himpunan dan
$\tau:\Cal PX\to[0,\infty]$ sembarang fungsi sedemikian sehingga $\tau\emptyset=0$.
Untuk $A\subseteq X$ tetapkan

$$\eqalign{\theta A
=\inf\{\sum_{j=0}\sp{\infty}\tau C_j:\langle C_j\rangle_{j\in\Bbb N}
\text{ adalah barisan }&\text{subhimpunan-subhimpunan }X\cr
&\text{sedemikian sehingga }
A\subseteq\bigcup_{j\in\Bbb N}C_j\}.\cr}$$

\noindent Tunjukkan bahwa $\theta$ adalah ukuran luar pada $X$. \Hint{anda
akan memerlukan 111F(b-ii) atau sesuatu yang setara.}

\spheader 113Ye Misalkan $X$ suatu himpunan dan $\theta_1$, $\theta_2$ dua
ukuran luar pada $X$. Tunjukkan bahwa $\theta_1\wedge\theta_2$, sebagaimana
diuraikan dalam 113Xb(iii), adalah ukuran luar yang diperoleh melalui proses
113Yd dari fungsional $\tau C=\min(\theta_1C,\theta_2C)$.

\spheader 113Yf Misalkan $X$ suatu himpunan dan
$\langle\theta_i\rangle_{i\in I}$ sembarang keluarga tak kosong yang terdiri
atas ukuran-ukuran luar pada $X$. Tetapkan $\tau C=\inf_{i\in I}\theta_iC$ untuk setiap
$C\subseteq X$. Tunjukkan bahwa ukuran luar yang diperoleh dari $\tau$ melalui
proses 113Yd adalah ukuran luar terbesar $\theta$ sedemikian sehingga
$\theta A\le\theta_iA$ setiap kali $A\subseteq X$ dan $i\in I$.

\spheader 113Yg Misalkan $X$ suatu himpunan dan $\phi:\Cal PX\to[0,\infty]$
suatu fungsional sedemikian sehingga

\inset{$\phi\emptyset=0$;

$\phi(A\cup B)\ge\phi A+\phi B$ untuk semua $A$, $B\subseteq X$ yang saling lepas;

jika $\sequencen{A_n}$ adalah barisan tak-menaik yang terdiri atas
subhimpunan-subhimpunan $X$ dan $\phi A_0<\infty$, maka
$\phi(\bigcap_{n\in\Bbb N}A_n)=\lim_{n\to\infty}\phi A_n$;

jika $\phi A=\infty$ dan $a\in\Bbb R$, maka terdapat $B\subseteq A$
sedemikian sehingga $a\le\phi B<\infty$.}

\noindent Tetapkan

\Centerline{$\Sigma=\{E:E\subseteq X,\,\phi(A\cap E)+\phi(A\setminus
E)=\phi A$ untuk setiap $A\subseteq X\}$.}

\noindent Tunjukkan bahwa $(X,\Sigma,\phi\restr\Sigma)$ adalah suatu ruang ukur.

\spheader 113Yh Misalkan $(X,\Sigma,\mu)$ suatu ruang ukur dan untuk
$A\subseteq X$ tetapkan
$\mu_*A=\sup\{\mu E:E\in\Sigma$, $E\subseteq A$, $\mu E<\infty\}$.
Tunjukkan bahwa $\mu_*:\Cal PX\to[0,\infty]$ memenuhi
syarat-syarat
113Yg, dan bahwa jika $\mu X<\infty$, maka ukuran yang didefinisikan dari $\mu_*$
dengan metode 113Yg memperluas $\mu$.

\spheader 113Yi Misalkan $X$ suatu himpunan dan $\Cal A$ suatu {\bf aljabar}
atas subhimpunan-subhimpunan $X$, yakni, suatu keluarga subhimpunan dari $X$ sedemikian sehingga

\inset{$\emptyset\in\Cal A$,

$X\setminus E\in\Cal A$ untuk setiap $E\in\Cal A$,

$E\cup F\in\Cal A$ setiap kali $E$, $F\in\Cal A$.}

\noindent Misalkan $\phi:\Cal A\to[0,\infty]$ suatu fungsi sedemikian sehingga

\inset{$\phi\emptyset=0$,

$\phi(E\cup F)=\phi E+\phi F$ setiap kali $E$, $F\in\Cal A$
dan $E\cap F=\emptyset$,

$\phi E=\lim_{n\to\infty}\phi E_n$ setiap kali
$\sequencen{E_n}$ adalah barisan tak-menurun di dalam $\Cal A$ yang gabungannya
$E$.}

\noindent Tunjukkan bahwa terdapat suatu ukuran $\mu$ pada $X$ yang memperluas $\phi$.
\Hint{tetapkan $\phi A=\infty$ untuk $A\in\Cal PX\setminus\Cal A$;
definisikan $\theta$ dari $\phi$ seperti dalam 113Yd, dan $\mu$ dari $\theta$.}

\spheader 113Yj (T.de Pauw) Misalkan $X$ suatu himpunan, $\Tau$ suatu
aljabar-$\sigma$ atas subhimpunan-subhimpunan $X$, dan $\theta$ suatu ukuran luar pada
$X$. Tetapkan
$\Sigma=\{E:E\in\Tau,\,\theta A=\theta(A\cap E)+\theta(A\setminus E)$
untuk setiap $A\in\Tau\}$. Tunjukkan bahwa $\Sigma$ adalah
aljabar-$\sigma$ atas subhimpunan-subhimpunan $X$ dan bahwa $\theta\restr\Sigma$ adalah suatu
ukuran.

\spheader 113Yk\dvAnew{2010} Misalkan $X$, $\tau:\Cal PX\to[0,\infty]$, dan
$\theta$ seperti dalam 113Yd; misalkan $\mu$ ukuran yang didefinisikan
dengan metode \Caratheodory\ dari $\theta$, dengan $\Sigma$ sebagai domain $\mu$.
Andaikan $E\subseteq X$ sedemikian sehingga
$\theta(C\cap E)+\theta(C\setminus E)\le\tau C$ setiap kali $C\subseteq X$
memenuhi $0<\tau C<\infty$. Tunjukkan bahwa $E\in\Sigma$.
}%end of exercises

\endnotes{
\Notesheader{113} Kita sedang menempuh tahap-tahap termudah yang dapat saya susun
menuju konstruksi suatu ruang ukur taktrivial, yakni ukuran Lebesgue
pada $\Bbb R$. Ada banyak konstruksi ukuran Lebesgue,
tetapi menurut saya metode \Caratheodory\ (113C) adalah metode yang tepat untuk memulai,
karena merupakan teknik tunggal yang paling ampuh dan serbaguna untuk
mengonstruksi ukuran. Tentu saja, metode ini abstrak -- metode ini menangani
sebarang ukuran luar pada sebarang himpunan; tetapi saya sungguh berpendapat bahwa
teori Lebesgue, yang terjalin dengan struktur ruang Euklides yang
kaya, lebih sulit daripada teori ukuran abstrak.
Setidaknya, di sini tersedia teorema berbobot yang dapat Anda
tekuni; menguasainya semestinya memuaskan sekaligus berguna. Saya
harus mengatakan bahwa menurut saya sungguh mengagumkan bahwa konstruksi langsung semacam itu
ternyata efektif. Dengan menelaah buktinya, mungkin ada gunanya
membedakan antara bagian `aljabar' atau `berhingga' ((a)-(c)) dan
bagian yang melibatkan barisan himpunan ((d)-(e)); bagian-bagian yang pertama merupakan
bukti 113Xh. Berbagai jenis ukuran luar muncul di seluruh
teori ukuran, dan saya menguraikan secara garis besar beberapa konstruksi yang relevan dalam
113X-113Y.
}%end of notes


\discrpage
